\DocumentMetadata{
  pdfversion=2.0,
  pdfstandard=ua-2,
  lang=en-US,
  tagging=on,
  tagging-setup={math/setup=mathml-SE,tabsorder=structure}
}
\documentclass[11pt,letterpaper]{article}
\usepackage[margin=0.68in,includefoot]{geometry}
\usepackage{newcomputermodern}
\usepackage{amsmath,amscd,mathtools}
\providecommand{\square}{\mdlgwhtsquare}
\usepackage[hidelinks]{hyperref}
\hypersetup{
  pdftitle={Analysis First-Year Exam, August 2018, Part 1},
  pdfauthor={University of Florida Department of Mathematics},
  pdfsubject={Graduate mathematics examination},
  pdfdisplaydoctitle=true,
  bookmarksopen=true
}
\setcounter{secnumdepth}{0}
\setlength{\parindent}{0pt}
\setlength{\parskip}{0.28em}
\setlength{\emergencystretch}{3em}
\setlength{\leftmargini}{2.15em}
\setlength{\leftmarginii}{2.65em}
\setlength{\leftmarginiii}{3em}
\renewcommand{\labelenumi}{\textbf{\theenumi.}}
\pagestyle{empty}
\raggedbottom
\allowdisplaybreaks
\newenvironment{examproblems}[1]{%
  \begin{enumerate}%
  \setcounter{enumi}{#1}%
  \setlength{\topsep}{0.35em}%
  \setlength{\partopsep}{0pt}%
  \setlength{\itemsep}{0.48em}%
  \setlength{\parsep}{0.18em}%
}{\end{enumerate}}
\newenvironment{examsubparts}{%
  \begin{enumerate}%
  \setlength{\topsep}{0.18em}%
  \setlength{\partopsep}{0pt}%
  \setlength{\itemsep}{0.16em}%
  \setlength{\parsep}{0.1em}%
}{\end{enumerate}}
\newenvironment{examinstructions}{%
  \par\begingroup\itshape
}{%
  \par\endgroup\vspace{0.18em}
}
\makeatletter
\renewcommand\section{\@startsection{section}{1}{0pt}%
  {-1.25ex plus -.25ex minus -.1ex}%
  {0.55ex plus .1ex}%
  {\normalfont\large\bfseries}}
\renewcommand{\@maketitle}{%
  \newpage\null\vskip -1.2em%
  {\footnotesize\bfseries
    \href{https://www.ufl.edu/}{University of Florida}, \href{https://math.ufl.edu/}{Department of Mathematics}\par}%
  \vskip 0.18em%
  \tagstructbegin{tag=Title}%
    {\LARGE\bfseries\href{https://gma.math.ufl.edu/exams/analysis-first-year/}{Analysis First-Year Exam}, August 2018, Part 1\par}%
  \tagstructend%
  \vskip 0.35em%
  \tagmcbegin{artifact}%
    \hrule height 1pt%
  \tagmcend%
  \vskip 0.55em%
}
\makeatother
\title{Analysis First-Year Exam, August 2018, Part 1}
\author{}
\date{}
\begin{document}
\maketitle
\thispagestyle{empty}
\begin{examinstructions}
Answer FOUR questions in detail. State carefully any results used without proof.
\end{examinstructions}
\begin{examproblems}{0}
\item
Let \((a_n)_{n=0}^{\infty}\) be a sequence of positive real numbers. Prove that 
\[
1/\limsup a_n=\liminf(1/a_n)
\]
 in the sense that if one side lies in \((0,\infty)\) then so does the other and the two sides are equal.
\item
Let \((a_n)_{n=0}^{\infty}\) be a sequence of real numbers and \(A=\{a_n:n\geq 0\}\) its set of values. For~\(p\) a real number, consider the two implications. One of these is true and the other false; give a proof for the true and a counterexample for the false.
\begin{examsubparts}
\item[\textnormal{(i)}]
if~\(p\) is a limit point of~\(A\) then~\(p\) is the limit of a subsequence of \((a_n)_{n=0}^{\infty}\);
\item[\textnormal{(ii)}]
if~\(p\) is the limit of a subsequence of \((a_n)_{n=0}^{\infty}\) then~\(p\) is a limit point of~\(A\).
\end{examsubparts}
\item
Let~\(X\) be a compact metric space and let \(f:X\to X\) be a continuous function. Prove that if~\(f\) has no fixed points, then there exists \(a>0\) such that \(d(x,f(x))\geq a\) whenever \(x\in X\). Show by example that the compact hypothesis cannot be dropped.
\item
This problem has two parts.
\begin{examsubparts}
\item[\textnormal{(i)}]
Prove that if~\(c\) is any real number then the following subset of the Euclidean plane is connected 
\[
\{(x,y):(x=\pm 1)\text{ or }(|x|<1\text{ and }y\geq c)\}.
\]
\item[\textnormal{(ii)}]
Let \(C_0\supseteq C_1\supseteq\cdots\) be a decreasing sequence of closed, connected subsets of a metric space. If the intersection \(\bigcap_{n\geq 0}C_n\) is nonempty, must it be connected? Proof or counterexample.
\end{examsubparts}
\item
Let \(f:(0,\infty)\to\mathbb{R}\) be differentiable and assume that \(f'(t)\to 0\) as the real number~\(t\) tends to infinity. Show that if \(f(n)\to\ell\) as the integer~\(n\) tends to infinity, then \(f(t)\to\ell\) as the real number~\(t\) tends to infinity.
\end{examproblems}
\end{document}
